%!TEX root = ../thesis.tex
%*******************************************************************************
%****************************** 4th Chapter **********************************
%*******************************************************************************
% Geo branch: area-data design (every Geo number regenerated by
% hull_geo_area_data_eval.py and hull_geo_area_data_models.py). Feature-branch
% text and numbers are unchanged from the previous version of this chapter.
\chapter{Empirical Validation of Geo-SMILE}
\label{chap:Empirical Validation}
% **************************** Define Graphics Path **************************
\ifpdf
    \graphicspath{{Chapter4/Figs/Raster/}{Chapter4/Figs/PDF/}{Chapter4/Figs/}}
\else
    \graphicspath{{Chapter4/Figs/Vector/}{Chapter4/Figs/}}
\fi

\section{Introduction}
\label{SEC:hull_intro}

This chapter presents the first independent empirical evaluation of the unified Geo-SMILE framework introduced in Chapter~3. The framework is evaluated in the context of next-month crime prediction in Kingston upon Hull, using monthly crime observations aggregated to Lower Layer Super Output Areas (LSOAs). The purpose of this experiment is to examine whether Geo-SMILE can generate faithful, stable, and spatially meaningful explanations when applied to a new prediction task, national context, spatial aggregation system, and data infrastructure that were not involved in the development of the framework.

Geo-SMILE extends the SMILE explainability framework of Aslansefat et al.\ \cite{aslansefat2021smile} by retaining its coalition-based perturbation procedure, distance-kernel weighting, and weighted linear surrogate, while introducing an explicit spatial explanation mechanism and distributional responses measured using the first Wasserstein distance. In the Hull experiment, the framework is applied through two complementary explanation branches designed to separate non-spatial feature sensitivity from geographic sensitivity.

The \emph{feature attribution branch} excludes coordinates and other explicit location variables and attributes model sensitivity across 30 non-spatial predictors. This branch is intended to explain how the predictive model responds to changes in the substantive explanatory variables without allowing geographic coordinates to compete directly with them within the attribution process.

The \emph{geo branch} explains the same predictive model from a geographic perspective. Its explanation players are the 168 Hull LSOAs, and perturbing an LSOA replaces all of the information the model receives from that location with a Hull-typical baseline. The branch therefore estimates, for each spatial unit, how sensitive the fitted model is to the data it takes from that location. Rather than treating latitude and longitude as conventional predictive features, this branch evaluates the model's dependence on each place as a distinct spatial component of the explanation.

The resulting Feature and Geo rankings are evaluated using both surrogate-based performance criteria and fixed expert-reference sets. These reference sets consist of the ten most important non-spatial predictors and the 20 most important LSOAs identified through a predefined expert-reference construction procedure. The construction, rationale, and standing of these reference sets are described in Section~\ref{sec:hull_ground_truth}. This provides an additional evaluation perspective beyond surrogate fidelity alone by examining whether Geo-SMILE identifies predictors and locations that are consistent with independently defined domain expectations.

Accordingly, this chapter focuses on the independent application and validation of Geo-SMILE within a crime-prediction setting. Unlike the benchmark experiments presented in the following chapter, the Hull study is not primarily designed as a direct comparison with an existing spatial explainability method. Instead, it investigates how the proposed framework behaves when transferred to a distinct real-world geospatial prediction problem and evaluates the resulting Feature and Geo explanations through complementary quantitative and expert-reference criteria.

\section{GeoSMILE for Next-Month Crime Prediction in Hull}
\label{sec:hull_first_independent_experiment}

The Kingston upon Hull study was designed as the first independent experiment conducted after the development of GeoSMILE. Its purpose was to determine whether the method could be transferred to a different city, national data infrastructure, spatial unit, and prediction task while retaining a strict separation between non-spatial feature importance and geographic sensitivity. The study was therefore treated as an independent application of the framework rather than as a direct replication of the development experiment.

The experiment addressed four questions. First, could GeoSMILE provide high-fidelity surrogate explanations for next-month crime prediction in a new urban setting? Second, were the resulting feature and geographic rankings reproducible across independently sampled perturbation sets? Third, did the highest-ranked factors exert a greater effect on the prediction model than random or bottom-ranked factors? Fourth, did the GeoSMILE rankings agree with the expert reference for important predictive features and LSOAs?

Only two active branches were evaluated. The \emph{Feature branch} used non-spatial predictors and explicitly excluded coordinates and location identifiers. The \emph{Geo branch} explained the same fitted model at the level of places, estimating one sensitivity coefficient for each of the 168 LSOAs by masking all of the model inputs belonging to that LSOA. Legacy Feature-A and Feature-B outputs were not included in this experiment.

\section{Study area, prediction target, and temporal design}
\label{sec:hull_data}

The observational unit was the LSOA-month. Police.uk street-crime records for Humberside Police were filtered to records assigned to Kingston upon Hull LSOAs \cite{policeuk}. A complete LSOA-by-month panel was constructed so that months without recorded offences were represented by zero counts rather than being omitted. The outcome variable was the total number of crimes recorded in the following month:

\begin{equation}
y_{i,t+1} = \text{crime count for LSOA } i \text{ in month } t+1.
\end{equation}

Temporal predictors were generated separately within each LSOA. These comprised one-, two-, and three-month crime lags; three- and six-month rolling means calculated from preceding months; and sine and cosine transformations of calendar month. Static contextual covariates were obtained from the 2021 Census and OpenStreetMap \cite{census2021,openstreetmap}. Official 2021 LSOA boundaries were used for mapping \cite{ons_lsoa2021}. Crime-event coordinates were aggregated to an LSOA location and transformed to British National Grid coordinates.

A chronological split was used to avoid leakage from future observations. Observations up to June 2025 formed the training period, July--December 2025 formed the validation period, and observations from January 2026 onward formed the test period. Training-set medians were used as perturbation baselines in both branches.

\subsection{Feature and geographic factors}
\label{sec:hull_factors}

The Feature branch contained 30 numeric, non-coordinate predictors. Coordinate fields, latitude, longitude, centroid variables, LSOA identifiers, and fields containing explicit location information were excluded by rule and checked at runtime. The Geo branch explained the same 30-predictor model, but its explanation players were the 168 individual LSOAs rather than predictor columns. Masking an LSOA replaced all 30 inputs of every LSOA-month observation belonging to that unit with their training medians. Each Geo coefficient therefore represented the sensitivity of the model to the data it receives from a particular LSOA.

\begin{table}[H]
\centering
\caption{Feature and Geo factors included in the Hull experiment.}
\label{tab:hull_factor_groups}
\small
\begin{tabularx}{\textwidth}{p{0.20\textwidth} p{0.09\textwidth} X}
\toprule
Factor group & Number & Variables or explanatory units \\
\midrule
Crime history & 5 &
One-, two-, and three-month crime lags; three-month rolling mean; six-month rolling mean. \\
Seasonality & 2 &
Sine and cosine transformations of calendar month. \\
Demographic structure & 3 &
Population density; percentage aged 15--34; percentage aged 65 and over. \\
Housing tenure & 3 &
Social renting; private renting; owner occupation. \\
Labour, health, and education & 5 &
Unemployment; economic inactivity; long-term sickness or disability; no qualifications; Level 4 or higher qualifications. \\
Built environment and exposure & 12 &
LSOA area; nightlife count and density; retail count and density; public-transport stop count and density; road length and density; intersection count and density; green-space percentage. \\
Geo factors & 168 &
One explanation player for each Hull LSOA. Masking a player replaces all 30 model inputs of that LSOA's observations with their training medians; the estimated coefficients are LSOA-specific. \\
\bottomrule
\end{tabularx}
\end{table}

\subsection{Selection of the predictive model by family}
\label{sec}

Geo-SMILE is model agnostic: the predictive model to which the explanation procedure is applied is not fixed by the method itself. This flexibility makes model selection an important part of the experimental design, because the resulting explanations describe the behaviour of a particular fitted predictor rather than the data-generating process independently of that predictor. The predictive model was therefore selected before the main Geo-SMILE analysis by comparing ten candidate models representing three broad model families.

\begin{description}

\item[Tree ensembles (XGBoost, LightGBM, CatBoost)]
These models represent gradient-boosted decision-tree approaches. They provide flexible non-linear response surfaces and can capture interactions between predictors without requiring them to be specified in advance.

\item[Kernel methods (Gaussian process, support vector regression)]
The Gaussian process and support vector regression provide a contrasting modelling family in which predictions are determined through kernel-based similarity rather than recursive partitioning of the predictor space.

\item[Neural architectures (ANN, CNN, MLP, Transformer, deep Keras network)]
Five neural architectures were included to test whether models based on learned distributed representations offered either a predictive or an explainability advantage over the tree and kernel alternatives.

\end{description}

The comparison was conducted in two stages. First, the candidates were compared on held-out predictive performance using the complete predictor set, including the spatial coordinates. Second, the fitted models were compared according to the fidelity with which the perturbation responses of the Feature and Geo branches could be represented by the surrogate models required by Geo-SMILE.

The purpose of the screening was not to identify a winner from small numerical differences. Instead, it was used to determine whether the candidate models formed clearly different performance groups and, within the strongest predictive group, whether any model also showed a clear advantage in surrogate fidelity.

\subsubsection{Stage 1: predictive adequacy}
\label{sec}

Table~\ref{tab:sel_pred} reports held-out predictive performance for the ten candidate models when all predictors, including the two spatial coordinates, are included.

The results show a distinct leading group. LightGBM achieves the highest predictive $R^2$ at $0.9045$, followed by XGBoost at $0.9000$, the Gaussian process at $0.8990$, and CatBoost at $0.8967$. The full range across these four models is only $0.0078$, indicating that their predictive performance is effectively very similar at the resolution relevant to the present comparison.

A second group is separated from this leading tier. ANN obtains an $R^2$ of $0.8557$, followed by CNN at $0.8368$, MLP at $0.8342$, and SVR at $0.8333$. The Transformer obtains $0.8079$, while the deep Keras network has the lowest predictive $R^2$ at $0.7837$.

The same separation is visible in the error metrics. The four leading models have MAE values between $4.6641$ and $4.8249$ and RMSE values between $6.5913$ and $6.8568$. The next-best model, ANN, has an MAE of $5.4571$ and an RMSE of $8.1037$.

Predictive performance therefore identifies four models---LightGBM, XGBoost, Gaussian process and CatBoost---as a common top tier. It does not, however, provide a convincing basis for selecting one of these four models over the others. The surrogate analysis therefore examines whether their explainability behaviour provides further separation.

\begin{table}[H]
\centering
\caption{Stage 1, held-out predictive performance for candidate models (all features, including coordinates)}
\label{tab:sel_pred}
\small
\begin{tabular}{llrrr}
\toprule
Family & Model & MAE & RMSE & $R^2$ \\
\midrule
\multirow{3}{*}{Tree ensemble}
 & LightGBM          & 4.6641 & 6.5913 & \textbf{0.9045} \\
 & XGBoost           & 4.6988 & 6.7477 & \textbf{0.9000} \\
 & CatBoost          & 4.8249 & 6.8568 & 0.8967 \\
\midrule
\multirow{2}{*}{Kernel}
 & Gaussian process  & 4.7461 & 6.7784 & 0.8990 \\
 & SVR               & 5.5525 & 8.7112 & 0.8333 \\
\midrule
\multirow{5}{*}{Neural}
 & ANN               & 5.4571 & 8.1037 & 0.8557 \\
 & CNN               & 5.3410 & 8.6173 & 0.8368 \\
 & MLP               & 5.1680 & 8.6864 & 0.8342 \\
 & Transformer       & 5.5237 & 9.3500 & 0.8079 \\
 & Keras (deep)      & 5.4330 & 9.9225 & 0.7837 \\
\bottomrule
\end{tabular}
\end{table}

\subsubsection{Stage 2: surrogate fidelity}
\label{sec}

Geo-SMILE requires a surrogate representation that produces one signed coefficient for each player. Consequently, the principal surrogate comparison focuses on linear models whose fitted parameters can be interpreted directly as player-level attribution coefficients. Non-linear surrogate models are also evaluated as diagnostic controls: they indicate whether additional response fidelity is available if the coefficient requirement is relaxed, but they cannot replace the linear surrogate in the downstream attribution analysis.

Tables~\ref{tab:fid_feat} and~\ref{tab:fid_geo} report surrogate fidelity for the Feature and Geo branches respectively, using $4{,}000$ perturbations in each case. For every candidate, the Geo branch explains the same fitted 30-predictor model as its Feature branch.

On the Feature branch, the principal linear estimators are extremely similar for most candidate predictive models. Weighted OLS, weighted Ridge, Bayesian ridge, ordinary linear regression and unweighted Ridge differ only in the third or fourth decimal place across the well-behaved models. For example, XGBoost obtains a weighted-Ridge $R^2$ of $0.9770$, compared with $0.9769$ for weighted OLS and Bayesian ridge and $0.9770$ for ordinary linear regression and unweighted Ridge. The same pattern occurs for LightGBM, CatBoost, the Gaussian process, the Transformer, ANN, SVR and CNN.

Among these models, Feature fidelity is highest for the Gaussian process ($R^2=0.9960$), followed by the Transformer ($0.9930$), LightGBM ($0.9832$), CatBoost ($0.9826$), XGBoost ($0.9770$), SVR ($0.9650$), CNN ($0.9275$), and ANN ($0.9265$) under weighted Ridge.

The two remaining neural models behave very differently. The deep Keras network produces a Feature weighted-Ridge $R^2$ of only $0.2037$, while the MLP produces $0.2286$. This is not a general failure of linear surrogates across the experiment, because the same surrogate obtains substantially higher fidelity for the other eight predictive models. These two models are therefore poor candidates for Feature-branch explanation under the Geo-SMILE linear-surrogate design.

The non-linear controls clarify the nature of the Feature-branch approximation. Histogram gradient boosting achieves higher Feature fidelity than weighted Ridge for every predictive model. For example, its $R^2$ is $0.9985$ for XGBoost, compared with $0.9770$ for weighted Ridge, and $0.9986$ for LightGBM, compared with $0.9832$. The effect is particularly large for the two poorly represented neural models: histogram gradient boosting raises fidelity to $0.6601$ for the deep Keras model and $0.6776$ for the MLP.

These improvements show that some additional non-linear structure remains in the perturbation response. They do not, however, make the non-linear surrogates suitable replacements for Ridge in Geo-SMILE, because the downstream procedure requires an explicit coefficient for every Feature or Geo player.

The Geo branch produces a different pattern. Linear surrogates dominate the non-linear alternatives across all ten predictive models, and they do so almost uniformly. Under weighted Ridge, Geo fidelity lies between $0.9454$ for the deep Keras network and $0.9485$ for SVR, with nine of the ten models between $0.9474$ and $0.9485$. This near-uniformity is a property of the Geo design: every masked observation receives the same baseline inputs, so each LSOA contributes a largely separable shift to the Wasserstein response, which a linear function of the LSOA masks represents well regardless of the underlying predictor.

Unweighted Ridge has the highest $R^2$ among the reported linear alternatives for each predictive model, but its advantage over weighted Ridge is numerically negligible. The largest difference is approximately $0.0006$, occurring for the MLP ($0.9480$ versus $0.9474$). For XGBoost, the corresponding values are $0.9486$ and $0.9484$; for CatBoost, $0.9482$ and $0.9479$; and for LightGBM, $0.9484$ and $0.9482$. Thus, the choice between the closely related linear estimators has essentially no effect on the model ordering.

By contrast, regularisation that induces stronger coefficient shrinkage performs poorly as the number of Geo players increases. The weak elastic net falls to $0.5591$ for XGBoost, $0.5600$ for LightGBM, $0.5605$ for CatBoost and $0.5731$ for the Gaussian process. Kernel ridge also performs substantially below the linear Ridge variants, with Geo $R^2$ values close to $0.62$ for every model.

The non-linear controls likewise fail to improve Geo fidelity. For XGBoost, histogram gradient boosting reaches $0.9047$ and random forest $0.8285$, both below the weighted-Ridge value of $0.9484$. The same pattern occurs throughout the table. This contrasts directly with the Feature branch, where histogram gradient boosting generally provides the highest raw fidelity.

Taken together, the two tables show that the principal linear estimators provide almost identical fidelity wherever a linear representation is adequate. \textbf{Weighted Ridge} is therefore retained as the common surrogate for the subsequent Geo-SMILE analysis. This choice should not be interpreted as a claim that it achieves the numerically highest $R^2$ in every individual cell; rather, its fidelity is effectively indistinguishable from the other leading linear estimators, while preserving the coefficient-based representation required by the downstream attribution analysis.

\begin{table}[H]
\centering
\caption{Feature branch fidelity results for candidate surrogates ($N=4{,}000$ perturbations)}
\label{tab:fid_feat}
\scriptsize
\setlength{\tabcolsep}{2.2pt}
\resizebox{\textwidth}{!}{%
\begin{tabular}{lrrrrrrrrrr}
\toprule
Surrogate & GP & Trans. & CatB. & LGBM & XGB & ANN & SVR & CNN & Keras & MLP \\
\midrule
Weighted OLS       & 0.9961 & 0.9930 & 0.9826 & 0.9832 & 0.9769 & 0.9265 & 0.9651 & 0.9274 & 0.2037 & 0.2286 \\
Weighted Ridge     & 0.9960 & 0.9930 & 0.9826 & 0.9832 & \textbf{0.9770} & 0.9265 & 0.9650 & 0.9275 & 0.2037 & 0.2286 \\
Bayesian ridge     & 0.9961 & 0.9930 & 0.9826 & 0.9832 & 0.9769 & 0.9265 & 0.9651 & 0.9274 & 0.2031 & 0.2277 \\
Linear regression  & 0.9961 & 0.9930 & 0.9827 & 0.9832 & 0.9770 & 0.9265 & 0.9651 & 0.9279 & 0.2042 & 0.2290 \\
Ridge (unweighted) & 0.9961 & 0.9930 & 0.9827 & 0.9832 & 0.9770 & 0.9265 & 0.9651 & 0.9279 & 0.2042 & 0.2290 \\
Elastic net (weak) & 0.9403 & 0.9269 & 0.9085 & 0.9383 & 0.9294 & 0.7949 & 0.8564 & 0.7874 & 0.1940 & 0.2205 \\
Random forest      & 0.9921 & 0.9885 & 0.9933 & 0.9971 & 0.9965 & 0.9160 & 0.9618 & 0.9210 & 0.8012 & 0.7910 \\
Hist.\ grad.\ boosting & 0.9976 & 0.9967 & 0.9985 & 0.9986 & 0.9985 & 0.9589 & 0.9798 & 0.9761 & 0.6601 & 0.6776 \\
Kernel ridge (RBF) & 0.9931 & 0.9920 & 0.9854 & 0.9860 & 0.9816 & 0.9369 & 0.9702 & 0.9498 & 0.3156 & 0.3448 \\
Tweedie GLM        & 0.5674 & 0.5879 & 0.5931 & 0.5752 & 0.5763 & 0.5028 & 0.5433 & 0.6126 & 0.1439 & 0.4596 \\
Poisson GLM        & 0.7852 & 0.7996 & 0.8020 & 0.7887 & 0.7882 & 0.6908 & 0.7455 & 0.8125 & 0.3293 & 0.7789 \\
\midrule
Unweighted OLS     & 0.9961 & 0.9930 & 0.9827 & 0.9832 & 0.9770 & 0.9265 & 0.9651 & 0.9279 & 0.2042 & 0.2290 \\
Ridge $\alpha=0$   & 0.9961 & 0.9930 & 0.9826 & 0.9832 & 0.9769 & 0.9265 & 0.9651 & 0.9274 & 0.2037 & 0.2286 \\
Shallow tree       & 0.9717 & 0.9179 & 0.8994 & 0.9877 & 0.9838 & 0.7394 & 0.9100 & 0.6935 & 0.2887 & 0.3332 \\
\midrule
Null, SD           & 0.0079 & 0.0075 & 0.0069 & 0.0069 & 0.0069 & 0.0067 & 0.0078 & 0.0078 & 0.0109 & 0.0082 \\
\bottomrule
\end{tabular}}
\end{table}


\begin{table}[H]
\centering
\caption{Geo branch fidelity results for candidate surrogates ($N=4{,}000$ perturbations; area-data masking of all inputs of the selected LSOAs)}
\label{tab:fid_geo}
\scriptsize
\setlength{\tabcolsep}{2.2pt}
\resizebox{\textwidth}{!}{%
\begin{tabular}{lrrrrrrrrrr}
\toprule
Surrogate & GP & Trans. & CatB. & LGBM & XGB & ANN & SVR & CNN & Keras & MLP \\
\midrule
Weighted OLS       & 0.9472 & 0.9475 & 0.9470 & 0.9473 & 0.9475 & 0.9473 & 0.9476 & 0.9472 & 0.9444 & 0.9462 \\
Weighted Ridge     & 0.9481 & 0.9483 & 0.9479 & 0.9482 & \textbf{0.9484} & 0.9483 & 0.9485 & 0.9481 & 0.9454 & 0.9474 \\
Bayesian ridge     & 0.9474 & 0.9477 & 0.9472 & 0.9475 & 0.9477 & 0.9476 & 0.9478 & 0.9475 & 0.9447 & 0.9465 \\
Linear regression  & 0.9482 & 0.9484 & 0.9481 & 0.9483 & 0.9485 & 0.9486 & 0.9486 & 0.9483 & 0.9456 & 0.9478 \\
Ridge (unweighted) & 0.9483 & 0.9485 & 0.9482 & 0.9484 & 0.9486 & 0.9487 & 0.9487 & 0.9484 & 0.9457 & 0.9480 \\
Elastic net (weak) & 0.5731 & 0.5430 & 0.5605 & 0.5600 & 0.5591 & 0.5876 & 0.5824 & 0.6150 & 0.4442 & 0.5773 \\
Random forest      & 0.8407 & 0.8169 & 0.8220 & 0.8254 & 0.8285 & 0.8558 & 0.8655 & 0.8957 & 0.7119 & 0.7976 \\
Hist.\ grad.\ boosting & 0.9120 & 0.8935 & 0.9064 & 0.9047 & 0.9047 & 0.9151 & 0.9200 & 0.9363 & 0.8390 & 0.8888 \\
Kernel ridge (RBF) & 0.6173 & 0.6164 & 0.6185 & 0.6184 & 0.6183 & 0.6185 & 0.6141 & 0.6122 & 0.6304 & 0.6246 \\
Tweedie GLM        & 0.1920 & 0.1908 & 0.1915 & 0.1923 & 0.1929 & 0.1911 & 0.1817 & 0.1945 & 0.1918 & 0.2179 \\
Poisson GLM        & 0.5116 & 0.4898 & 0.4998 & 0.5037 & 0.5031 & 0.5318 & 0.4931 & 0.5858 & 0.4398 & 0.5946 \\
\midrule
Unweighted OLS     & 0.9482 & 0.9484 & 0.9481 & 0.9483 & 0.9485 & 0.9486 & 0.9486 & 0.9483 & 0.9456 & 0.9478 \\
Ridge $\alpha=0$   & 0.9472 & 0.9475 & 0.9470 & 0.9473 & 0.9475 & 0.9473 & 0.9476 & 0.9472 & 0.9444 & 0.9462 \\
Shallow tree       & 0.7508 & 0.7243 & 0.7258 & 0.7331 & 0.7332 & 0.7999 & 0.8177 & 0.8533 & 0.4661 & 0.7231 \\
\midrule
Null, SD           & 0.0161 & 0.0185 & 0.0163 & 0.0174 & 0.0174 & 0.0159 & 0.0174 & 0.0170 & 0.0164 & 0.0149 \\
\bottomrule
\end{tabular}}
\end{table}


\subsubsection{Comparison under the selected surrogate}
\label{sec}

Table~\ref{tab:ridge_fidelity_full} brings the two branches together using weighted Ridge for all ten predictive models.

The results reinforce the distinction between predictive-model selection and surrogate fidelity. A model can be highly predictable by the surrogate without being among the strongest predictive models, and strong fidelity on one branch does not guarantee equally strong behaviour on the other.

The Gaussian process has the highest Feature fidelity in the complete comparison, with $R^2=0.9960$, MAE $=0.1203$, and weighted RMSE $=0.1588$. The Transformer is also very strong on Feature, with $R^2=0.9930$, despite its lower predictive $R^2$ of $0.8079$ in Stage~1. Conversely, the deep Keras network and MLP show the most pronounced branch imbalance. Their Geo fidelities are $0.9454$ and $0.9474$, respectively, but their Feature fidelities fall to $0.2037$ and $0.2286$.

The four models identified as the predictive top tier remain strong under the surrogate analysis. Their results are:

\begin{itemize}

\item LightGBM: predictive $R^2=0.9045$, Feature fidelity $R^2=0.9832$, and Geo fidelity $R^2=0.9482$;

\item XGBoost: predictive $R^2=0.9000$, Feature fidelity $R^2=0.9770$, and Geo fidelity $R^2=0.9484$;

\item Gaussian process: predictive $R^2=0.8990$, Feature fidelity $R^2=0.9960$, and Geo fidelity $R^2=0.9481$;

\item CatBoost: predictive $R^2=0.8967$, Feature fidelity $R^2=0.9826$, and Geo fidelity $R^2=0.9479$.

\end{itemize}

No member of this group dominates the others across the three criteria. LightGBM has the highest predictive $R^2$, the Gaussian process has the highest Feature surrogate fidelity, and XGBoost has the highest Geo weighted-Ridge fidelity.

The differences are small. Across the four models, predictive $R^2$ spans only $0.0078$, Feature fidelity spans $0.0190$, and Geo fidelity spans only $0.0005$. The tables therefore support the conclusion that these four models form a high-performing group, but they do not support declaring any one of them the unique empirical winner.

\begin{table}[H]
\centering
\caption{Weighted Ridge surrogate fidelity, all ten models, both branches ($N=4{,}000$ perturbations each).}
\label{tab:ridge_fidelity_full}
\scriptsize
\setlength{\tabcolsep}{2.2pt}
\resizebox{\textwidth}{!}{%
\begin{tabular}{llrrrrrrrrrr}
\toprule
& & \multicolumn{5}{c}{Feature} & \multicolumn{5}{c}{Geo} \\
\cmidrule(lr){3-7}\cmidrule(lr){8-12}
Family & Model & $R^2$ & MAE & WMAE & RMSE & WRMSE & $R^2$ & MAE & WMAE & RMSE & WRMSE \\
\midrule
\multirow{3}{*}{Tree ensemble}
 & XGBoost   & 0.9770 & 0.3591 & 0.3590 & 0.4094 & 0.4095 & 0.9484 & 0.7849 & 0.8007 & 1.3550 & 1.4001 \\
 & LightGBM  & 0.9832 & 0.3262 & 0.3261 & 0.3742 & 0.3739 & 0.9482 & 0.7950 & 0.8109 & 1.3781 & 1.4239 \\
 & CatBoost  & 0.9826 & 0.1721 & 0.1737 & 0.2288 & 0.2313 & 0.9479 & 0.7999 & 0.8148 & 1.3779 & 1.4227 \\
\midrule
\multirow{2}{*}{Kernel}
 & Gaussian process & 0.9960 & 0.1203 & 0.1202 & 0.1576 & 0.1588 & 0.9481 & 0.8585 & 0.8740 & 1.5213 & 1.5720 \\
 & SVR & 0.9650 & 0.2073 & 0.2056 & 0.2667 & 0.2653 & 0.9485 & 0.7712 & 0.7821 & 1.5003 & 1.5506 \\
\midrule
\multirow{5}{*}{Neural}
 & ANN         & 0.9265 & 0.3427 & 0.3426 & 0.4452 & 0.4466 & 0.9483 & 0.9416 & 0.9534 & 1.7660 & 1.8215 \\
 & CNN         & 0.9275 & 0.2836 & 0.2852 & 0.3683 & 0.3718 & 0.9481 & 1.1968 & 1.2119 & 2.4303 & 2.5130 \\
 & MLP & 0.2286 & 732.6481 & 730.8138 & 1392.2525 & 1391.9188 & 0.9474 & 1.1014 & 1.1110 & 1.9185 & 1.9729 \\
 & Transformer & 0.9930 & 0.1468 & 0.1461 & 0.1882 & 0.1862 & 0.9483 & 0.7134 & 0.7287 & 1.2309 & 1.2731 \\
 & Keras (deep) & 0.2037 & 8.0351 & 8.0713 & 16.6323 & 16.7122 & 0.9454 & 0.5616 & 0.5726 & 0.8164 & 0.8399 \\
\bottomrule
\end{tabular}}
\end{table}


\subsubsection{Selected model}
\label{sec}

The screening therefore leads to a deliberately limited conclusion. The evidence identifies a top tier consisting of LightGBM, XGBoost, CatBoost and the Gaussian process, rather than a single best model.

Within this group, each model has a different numerical advantage. LightGBM gives the strongest held-out prediction, the Gaussian process gives the strongest Feature-branch Ridge fidelity, and XGBoost gives the strongest Geo-branch Ridge fidelity ($0.9484$), while remaining close to the leaders on the other two criteria, with a predictive $R^2$ of $0.9000$ and Feature fidelity of $0.9770$.

Accordingly, XGBoost is adopted as the reference predictive model for the remainder of this chapter, but its selection is not presented as the outcome of a numerical victory over LightGBM, CatBoost or the Gaussian process. Rather, the quantitative screening establishes that XGBoost belongs to the same high-performing group and that using it does not require accepting a material loss in either predictive performance or surrogate fidelity.

This distinction is important for the interpretation of the subsequent Geo-SMILE results. The purpose of the screening is to show that the explanations are being generated for a predictor that is both competitively accurate and well represented by the selected surrogate in both branches. The results do not imply that XGBoost is uniquely required by Geo-SMILE, and the numerical evidence presented here would also support LightGBM, CatBoost or the Gaussian process as defensible members of the same top-performing group.



\subsection{Predictive and explanation models}
\label{sec:hull_models}

After the model-family and surrogate screening, XGBoost was retained as the reference predictive model and weighted Ridge as the common Geo-SMILE surrogate. A single XGBoost count model $f$ was fitted on the 30 non-spatial predictors, using a Poisson objective, 500 trees, a learning rate of 0.04, maximum depth of five, minimum child weight of three, subsampling and column-subsampling rates of 0.85, an $L_2$ regularisation parameter of 2.0, and random seed 42. Both branches explain this one model: the Feature branch partitions its sensitivity across predictors, and the Geo branch partitions it across places. The two branches are therefore complementary views of one predictor rather than additive components of a combined prediction.

For the Feature branch, binary masks were generated. A value of one retained a feature, whereas a value of zero replaced that feature with its training-set median. For perturbation $m$, the response was the first Wasserstein distance between the original and perturbed test prediction distributions:
\begin{equation}
d_m^{(F)} =
W_1\left[
f(X),
f\left(\widetilde{X}^{(m)}\right)
\right].
\label{eq:hull_feature_response}
\end{equation}

For the Geo branch, each perturbation selected between five and ten LSOAs without replacement and replaced \emph{all} 30 inputs of every affected LSOA-month observation with the corresponding training medians. The response was calculated only over the affected observations:
\begin{equation}
d_m^{(G)} =
W_1\left[
f(X_{A_m}),
f\left(\widetilde{X}^{(m)}_{A_m}\right)
\right],
\label{eq:hull_geo_response}
\end{equation}
where $A_m$ denotes the LSOA-month observations of the LSOAs selected in perturbation $m$. Because every masked observation receives the same baseline inputs, its perturbed prediction equals a single \emph{typical-area} prediction; an LSOA therefore contributes more to $d_m^{(G)}$ the further its own predictions lie from that typical value. The Geo coefficient of an LSOA consequently measures how strongly the fitted model depends on the data it receives from that place.

The final publication design uses $N=4{,}000$ perturbations in \emph{both} branches. Perturbation masks and responses were fitted with kernel-weighted Ridge surrogates with $\alpha=1$. RBF, Laplace and inverse kernels were compared on the validation split, and the mask pool was divided into 60\% surrogate training, 20\% validation and 20\% testing. Feature importance was ranked by the absolute coefficient, with the signed coefficient retained only to describe the direction of the model response under the baseline-ablation experiment. Geo importance was ranked by the \emph{signed} coefficient. In the Geo surrogate, the coefficient of an LSOA is its contribution relative to the average LSOA contribution absorbed by the intercept; a strongly negative coefficient therefore denotes an LSOA whose data the model barely depends on, not a strong effect of opposite direction, and taking absolute values would wrongly promote such LSOAs to the top of the ranking.

The move to a common $N=4{,}000$ reporting budget replaces the earlier asymmetric publication design of 1,024 Feature masks and 2,000 Geo perturbations. It also aligns the final Hull reporting design with the $N=4{,}000$ model/surrogate screening reported above.

\begin{table}[H]
\centering
\caption{Final research design for the Feature and Geo branches.}
\label{tab:hull_research_design}
\small
\begin{tabularx}{\textwidth}{p{0.24\textwidth} X X}
\toprule
Design element & Feature branch & Geo branch \\
\midrule
Purpose & Explain non-spatial predictive factors & Explain LSOA-specific sensitivity to each place's data \\
Predictive inputs & 30 non-coordinate variables & The same 30 variables (same fitted model) \\
Explanation players & 30 features & 168 LSOAs \\
Predictive model & XGBoost Poisson count model & The same XGBoost model \\
Reporting perturbations & 4,000 binary masks & 4,000 LSOA perturbations \\
Perturbed units & Any binary subset of features & Random sets of 5--10 LSOAs \\
Masking baseline & Training median of each feature & Training median of all 30 inputs, applied to every observation of the selected LSOAs \\
Perturbation response & Global $W_1$ between original and perturbed prediction distributions & $W_1$ over predictions for affected LSOA-month observations \\
Kernel distance & Fraction of features masked & Fraction of LSOAs masked \\
Kernel candidates & RBF, Laplace, inverse & RBF, Laplace, inverse \\
Surrogate & Weighted Ridge, $\alpha=1$ & Weighted Ridge, $\alpha=1$ \\
Surrogate mask split & 60\% training, 20\% validation, 20\% testing & 60\% training, 20\% validation, 20\% testing \\
Importance score & Absolute feature coefficient & Signed LSOA coefficient \\
Expert-reference comparison & Top 10 of 30 & Top 20 of 168 \\
Faithfulness experiment & Top 10, bottom 10, and 1,000 random sets of 10 & Top 20, bottom 20, and 1,000 random sets of 20 \\
\bottomrule
\end{tabularx}
\end{table}

\subsection{Results}
\label{sec:hull_results}

All quantitative evaluation results in the remainder of this section use the $N=4{,}000$ metric-suite run for both branches. The local maps supplied with the final result package are likewise $N=4{,}000$ outputs; the six mapped Feature drivers are the same six leading variables identified by the metric-suite ranking, although small coefficient differences can occur between separately exported perturbation runs.

\subsubsection{Global Feature importance}

Figure~\ref{fig:hull_global_feature_importance} presents the global Feature-SMILE ranking at $N=4{,}000$. Importance is defined by the absolute weighted-Ridge coefficient.

\begin{figure}[H]
\centering
\includegraphics[width=0.8\textwidth]{Chapter4/Figs/top_global_drivers_N4000.png}
\caption{Leading global Feature-SMILE drivers of next-month crime prediction in Hull}
\label{fig:hull_global_feature_importance}
\end{figure}

The ranking remains strongly dominated by recent crime history. The six-month rolling crime mean has the largest coefficient magnitude ($5.1425$), followed by the three-month rolling crime mean ($1.2502$). The one-month crime lag is third ($0.3940$), followed by unemployment ($0.3608$), long-term sickness or disability ($0.2966$), and owner occupation ($0.2179$). Social renting ($0.1206$), road length ($0.1032$), no qualifications ($0.0952$), and private renting ($0.0937$) complete the top ten. All ten coefficients are negative in the fitted surrogate, but this sign should be interpreted only as the direction of the baseline-ablation response; importance is assessed by magnitude.

\subsubsection{Geographic sensitivity}
\label{subsubsec:geographic_sensitivity}

A key contribution of Geo-SMILE is the explicit separation of
\emph{geographic sensitivity} from the effects of individual predictors.
While the Feature branch explains which variables the fitted model relies on,
the Geo branch identifies \emph{where} the same model relies most on the data
it receives. At $N=4{,}000$, the resulting Geo attributions show substantial
spatial heterogeneity across Hull
(Figure~\ref{fig:hull_geo_sensitivity}).

\begin{figure}[H]
    \centering
    \includegraphics[width=\textwidth]{Chapter4/Figs/geo_area_data_map_N4000.pdf}
    \caption{Geographic sensitivity identified by the Geo-SMILE Geo branch for
    Hull. (a) Signed Geo-SMILE coefficient; (b) prediction shift, in crimes
    per month, when all of an LSOA's data are replaced by the Hull training
    median, with the expert-reference top 20 outlined.}
    \label{fig:hull_geo_sensitivity}
\end{figure}

Panel (a) presents the signed Geo-SMILE coefficients, which define the Geo
ranking. Positive values identify LSOAs whose data the model depends on more
than the average LSOA; negative values identify LSOAs whose predictions stay
close to the typical-area prediction when their data are removed, and
therefore places on which the model depends comparatively little. The
coefficient ordering agrees closely with the directly computed per-LSOA
prediction shift in panel (b) (Spearman $\rho=0.982$), which confirms that the
surrogate ranking reflects the model's actual dependence on each place.

The distribution is highly skewed. \texttt{E01033104} (Kingston upon Hull 029E)
has by far the largest coefficient ($22.33$), corresponding to a prediction
shift of $176.7$ crimes per month when its data are masked. It is followed by
\texttt{E01012889} ($7.18$), \texttt{E01012871} ($4.73$), \texttt{E01012854}
($4.27$), and \texttt{E01012852} ($4.13$). The highest-sensitivity areas form a
contiguous concentration along the southern and central riverside, with
further isolated high-sensitivity units in the north and east, whereas most of
the northern, western and eastern suburbs show below-average sensitivity. The
median per-LSOA prediction shift is $6.4$ crimes per month, and the 90th
percentile is $22.0$.

Together, the two panels distinguish the surrogate ranking from the
magnitude of the underlying dependence, and they provide a spatially explicit
account of where the fitted model relies on local information. As with the
other Geo-SMILE attributions, these coefficients describe model sensitivity
and should not be interpreted as causal effects of geographic location.


\subsubsection{Local Feature importance}

Local Feature-SMILE coefficients were calculated by LSOA to examine whether the strongest global drivers are spatially homogeneous. At $N=4{,}000$, the six mapped drivers are the six-month rolling crime mean, three-month rolling crime mean, one-month crime lag, unemployment, long-term sickness or disability, and owner occupation. The important correction relative to the earlier publication figure is that \texttt{crime\_lag\_1}, not \texttt{crime\_lag\_2}, belongs in the six-panel figure.

\begin{figure}[H]
\centering
\includegraphics[width=\textwidth]{Chapter4/Figs/Top_drivers_global_importance_N4000.png}
\caption{Spatial expression of the six leading global Feature-SMILE drivers at $N=4{,}000$. Panels show signed local coefficients for six-month rolling crime mean, three-month rolling crime mean, one-month crime lag, unemployment, long-term sickness or disability, and owner occupation. Red and blue denote opposite directions of local model response; stronger intensity denotes larger coefficient magnitude.}
\label{fig:hull_top_feature_drivers}
\end{figure}

The two rolling crime measures remain dominant but vary markedly across neighbourhoods. The local maps also show that the one-month lag and the socioeconomic predictors do not act uniformly across the city. The maps should therefore be read as spatial variation in the fitted model's local sensitivity rather than as neighbourhood-level causal effects.

\subsubsection{Surrogate fidelity and perturbation budget}
\label{sec:hull_fidelity_n}

The final publication design standardises both branches at $N=4{,}000$. The perturbation sweep shows why this is defensible without claiming that every diagnostic is maximised at exactly the same $N$. Feature fidelity is already saturated throughout the tested range, remaining between approximately $0.975$ and $0.977$. Geo fidelity improves materially between the smaller budgets and reaches its highest held-out value in the sweep at $N=4{,}000$, $R^2=0.9484$. At the same budget, both branches have top-$k$ stability Jaccard equal to $1.000$, while rank consistency rises to $0.831$ for Feature and $0.963$ for Geo.

\begin{table}[H]
\centering
\caption{Held-out Geo-SMILE weighted-Ridge fidelity by perturbation count}
\label{tab:hull_fidelity_by_n}
\small
\begin{tabular}{rrrrr}
\toprule
$N$ & Feature $R^2$ & Feature weighted $R^2$ & Geo $R^2$ & Geo weighted $R^2$ \\
\midrule
1,000 & 0.9764 & 0.9765 & 0.9309 & 0.9297 \\
2,000 & 0.9774 & 0.9773 & 0.9231 & 0.9365 \\
3,000 & 0.9768 & 0.9770 & 0.9440 & 0.9414 \\
\textbf{4,000} & \textbf{0.9770} & \textbf{0.9771} & \textbf{0.9484} & \textbf{0.9461} \\
5,000 & 0.9766 & 0.9770 & 0.9418 & 0.9427 \\
7,500 & 0.9771 & 0.9772 & 0.9421 & 0.9428 \\
10,000 & 0.9753 & 0.9753 & 0.9459 & 0.9456 \\
\bottomrule
\end{tabular}
\end{table}



The common $N=4{,}000$ budget is therefore used for the maps and for the quantitative evaluation below. This choice should be interpreted as a reporting standardisation supported by the sweep, not as evidence that $N=4{,}000$ uniquely optimises every criterion.

\subsection{Evaluation metrics}
\label{sec:hull_quant_results}

The quantitative evaluation follows the five attribution-evaluation headings used for SMILE: ATT fidelity, ATT accuracy, faithfulness, ATT stability and ATT consistency.

\paragraph{Provenance of the expert reference}
\label{sec:hull_ground_truth}

The Feature reference is the fixed top ten of the 30 ranked predictors, and the Geo reference is the fixed top 20 of the 168 Hull LSOAs via expert reference. Only the Geo-SMILE ranking changes when the perturbation design is rerun.

\begin{table}[H]
\centering
\caption{Composition of the expert-reference sets used in the Hull experiment.}
\label{tab:hull_ground_truth_design}
\small
\begin{tabularx}{\textwidth}{p{0.24\textwidth} p{0.09\textwidth} X}
\toprule
Reference component & Number & Members and rationale \\
\midrule
Feature factors, crime and temporal & 4 &
\texttt{crime\_lag\_1}, \texttt{crime\_rolling\_mean\_3}, \texttt{crime\_lag\_2}, and \texttt{crime\_rolling\_mean\_6}. \\
Feature factors, economic, labour and tenure & 4 &
\texttt{unemployed\_pct\_age16plus}, \texttt{economically\_inactive\_pct}, \texttt{social\_rented\_pct}, and \texttt{owner\_occupied\_pct}. \\
Feature factors, built environment and activity & 2 &
\texttt{nightlife\_per\_km2} and \texttt{retail\_per\_km2}. \\
Geo reference LSOAs & 20 &
The $N=4{,}000$ metric-suite Geo top 20 overlaps the reference in 9 units:
\texttt{E01012763},
\texttt{E01012805},
\texttt{E01012852},
\texttt{E01012854},
\texttt{E01012855},
\texttt{E01012871},
\texttt{E01012889},
\texttt{E01012891},
\texttt{E01033104}. \\
\bottomrule
\end{tabularx}
\end{table}

The rating remains a single-rater expert reference. The resulting agreement statistics therefore quantify correspondence with that reference rather than independent validation of causal importance.


\subsubsection{ATT fidelity}
\label{sec:hull_att_fidelity}

At the common $N=4{,}000$ design, the Feature weighted-Ridge surrogate obtains held-out $R^2=0.9770$ and weighted $R^2=0.9771$. The Geo surrogate obtains held-out $R^2=0.9484$ and weighted $R^2=0.9461$. Thus both branches are well approximated by the selected linear surrogate.

\begin{table}[H]
\centering
\caption{ATT fidelity at the final $N=4{,}000$ design (seed 42).}
\label{tab:hull_fidelity_selected}
\small
\resizebox{\textwidth}{!}{%
\begin{tabular}{lcccccc}
\toprule
Branch &
Held-out $R^2$ &
Weighted $R^2$ &
MAE &
Weighted MAE &
RMSE &
Weighted RMSE \\
\midrule
Feature & 0.9770 & 0.9771 & 0.3591 & 0.3590 & 0.4094 & 0.4095 \\
Geo     & 0.9484 & 0.9461 & 0.7849 & 0.8007 & 1.3550 & 1.4001 \\
\bottomrule
\end{tabular}}
\end{table}


\subsubsection{ATT accuracy}
\label{sec:hull_att_accuracy}

At $N=4{,}000$, the Feature branch recovers six of the ten reference predictors. Precision and recall at $k=10$ are both $0.60$, Jaccard similarity is $0.4286$, ATT ACC is $0.7333$, ATT F1 is $0.6000$, and ATT AUROC is $0.8300$. The six shared predictors are \nolinkurl{crime_rolling_mean_6}, \nolinkurl{crime_rolling_mean_3}, \nolinkurl{crime_lag_1}, \nolinkurl{unemployed_pct_age16plus}, \nolinkurl{owner_occupied_pct}, and \nolinkurl{social_rented_pct}. The four reference items outside the model top ten are \nolinkurl{crime_lag_2}, \nolinkurl{economically_inactive_pct}, \nolinkurl{nightlife_per_km2}, and \nolinkurl{retail_per_km2}; the four model-only top-ten items are \nolinkurl{long_term_sick_disabled_pct}, \nolinkurl{road_length_km}, \nolinkurl{no_qualifications_pct}, and \nolinkurl{private_rented_pct}.

The Geo branch recovers 9 of the 20 reference LSOAs. Precision and recall are both $0.45$, Jaccard similarity is $0.2903$, ATT ACC is $0.8690$, ATT F1 is $0.4500$, and ATT AUROC is $0.7250$. Agreement is concentrated at the top of the ranking: seven of the eight highest-ranked LSOAs, including all of the five most sensitive, are reference members. Several reference-only LSOAs rank just outside the top 20 (for example \texttt{E01012757} at 22nd and \texttt{E01012868} at 23rd).

\begin{table}[H]
\centering
\caption{Primary ATT accuracy at the final $N=4{,}000$ design.}
\label{tab:hull_att_accuracy_primary}
\small
\begin{tabular}{lcc}
\toprule
Metric & Feature & Geo \\
\midrule
$n$             & 30     & 168 \\
$k$             & 10     & 20 \\
Intersection    & 6      & 9 \\
Accuracy        & 0.6000 & 0.4500 \\
ATT ACC         & 0.7333 & 0.8690 \\
ATT F1          & 0.6000 & 0.4500 \\
ATT AUROC       & 0.8300 & 0.7250 \\
Precision@$k$   & 0.6000 & 0.4500 \\
Recall@$k$      & 0.6000 & 0.4500 \\
Jaccard@$k$     & 0.4286 & 0.2903 \\
\bottomrule
\end{tabular}
\end{table}

% The previous Geo expert-comparison map encoded the coordinates-only ranking.
% Regenerate that map from the area-data N=4,000 Geo top 20 before restoring the figure.

\subsubsection{ATT Faithfulness}
\label{sec:hull_faithfulness}

ATT Faithfulness was evaluated across the candidate predictive models using
the performance--attribution correspondence defined in
Section~\ref{subsec:faithfulness}. For each branch, the held-out predictive
$R^2$ associated with each model was paired with the attribution quality
obtained from that model's Geo-SMILE explanation. Because both branches now
explain the same fitted 30-predictor model, the Geo pairing uses that model's
predictive $R^2$. ATT-AUC was treated as the primary attribution-quality
measure, while top-$k$ overlap with the expert reference was evaluated as a
secondary measure.

Table~\ref{tab:hull_faithfulness_correlation} reports the resulting
correlations. No statistically reliable positive association was detected in
either branch.

For the Feature branch, predictive $R^2$ was essentially unrelated to
ATT-AUC ($r=0.10$, two-sided $p=0.79$) and showed a weak negative association
with top-10 overlap ($r=-0.12$, $p=0.74$). The corresponding Spearman
correlations were also close to zero. Predictive competence therefore did not
systematically track attribution agreement across the ten Feature models.

The Geo branch, now evaluated for all ten models, showed no positive
tendency. The association between predictive $R^2$ and ATT-AUC was
$r=-0.00$ ($p=0.99$), while the association with top-20 overlap was
$r=-0.48$ ($p=0.16$). Neither relationship was statistically significant, and
the negative overlap coefficient runs against the direction expected by the
faithfulness hypothesis.

\begin{table}[H]
\centering
\caption{ATT Faithfulness in the Hull experiment, measured as the association
between predictive $R^2$ and attribution quality across candidate models.}
\label{tab:hull_faithfulness_correlation}
\small
\begin{tabular}{llrrrr}
\toprule
Branch & Attribution measure & $M$ & Pearson $r$ &
Two-sided $p$ & Spearman $\rho$ \\
\midrule
Feature & Top-10 reference overlap & 10 & -0.12 & 0.74 & -0.09 \\
Feature & ATT-AUC                  & 10 &  0.10 & 0.79 & -0.09 \\
Geo     & Top-20 reference overlap & 10 & -0.48 & 0.16 & -0.39 \\
Geo     & ATT-AUC                  & 10 & -0.00 & 0.99 & -0.10 \\
\bottomrule
\end{tabular}
\end{table}

The model-level results explain why the aggregate correlations are weak.
Within the Feature branch, the Transformer combines the lowest predictive
$R^2$ in the model set with one of the highest ATT-AUC values, whereas the
deep-learning and MLP models occupy an intermediate predictive range but have
the lowest ATT-AUC values. Predictive ordering therefore does not translate
into attribution-quality ordering.

For the Geo branch, attribution quality varies within a narrow band. The
boosting models and the Gaussian process, which have the highest predictive
$R^2$ (0.896--0.899), obtain ATT-AUC values between $0.725$ and $0.748$ and
overlaps of 9--10 of 20. The MLP, with a lower predictive $R^2$ of $0.853$,
has the highest ATT-AUC ($0.867$) and overlap (11/20), whereas the
Transformer ($R^2=0.873$) and SVR ($0.863$) have the lowest ATT-AUC values
($0.693$ and $0.694$). The two weakest predictors in this run, the deep Keras
network ($R^2=0.666$) and the CNN ($0.635$), obtain overlaps of 11 and 10 of
20. The predictive $R^2$ values of these two Keras models were lower in this
run than in Stage~1 because Keras training is not fully deterministic under a
fixed seed; the Geo pairing uses the $R^2$ of the models whose explanations
were actually evaluated.

The within-family correlations are not used as primary evidence because the
model counts are too small for stable inference and, for the boosting models,
predictive $R^2$ differs only minimally between models.

These results do not demonstrate ATT Faithfulness under the
performance--attribution correspondence criterion. They instead indicate that,
within the present Hull model set, attribution agreement with the expert
reference varies for reasons not captured by predictive $R^2$ alone. The
limited number of paired models also restricts statistical power: with only
ten observations, only comparatively strong correlations would be detectable
at conventional significance levels. The result is therefore interpreted as an
absence of positive evidence for this criterion, rather than as evidence that
Geo-SMILE explanations are intrinsically unfaithful.


\subsubsection{ATT stability}
\label{sec:hull_att_stability}

ATT stability evaluates whether the Geo-SMILE explanation is preserved after a
small, controlled modification of the explanation input. Stability was tested
by expanding the Hull explanation reference with one real neighbouring LSOA in
each cardinal direction (north, east, south, and west), one addition at a time.
The fitted XGBoost model, Hull observations, surrogate specification,
perturbation budget ($N=4{,}000$), and SMILE seed (42) were held fixed.
Consequently, differences between the baseline and expanded explanations can be
attributed to the controlled spatial addition rather than to a different
predictive model or perturbation sample.

For each addition, stability was assessed using two complementary quantities.
First, Jaccard similarity measured preservation of the highest-ranked
attribution set ($k=10$ for Feature and $k=20$ for Geo). Second, Spearman rank
correlation measured agreement over the complete coefficient vector (absolute
coefficients for Feature, signed coefficients for Geo). Prediction Spearman
correlation was additionally examined to verify that the ordering of the
original Hull predictions remained unchanged.

\begin{table}[H]
\centering
\caption{ATT stability at $N=4{,}000$ (seed 42) under one-LSOA additions in
the four cardinal directions.}
\label{tab:hull_att_stability}
\scriptsize
\resizebox{\textwidth}{!}{%
\begin{tabular}{lcc|cc|cc|cc}
\toprule
Branch &
North Jac. & North $\rho$ &
East Jac. & East $\rho$ &
South Jac. & South $\rho$ &
West Jac. & West $\rho$ \\
\midrule
Feature &
1.000 & 0.9996 &
1.000 & 0.9996 &
1.000 & 1.0000 &
1.000 & 0.9991 \\
Geo &
1.000 & 0.9998 &
1.000 & 0.9998 &
1.000 & 0.9998 &
1.000 & 0.9998 \\
\bottomrule
\end{tabular}}
\end{table}

The highest-ranked attribution sets were unchanged in every cardinal-direction
trial. Feature top-10 and Geo top-20 Jaccard similarity were both $1.000$ for
all four additions. Agreement over the complete attribution vectors was also
extremely high. Mean Spearman correlation was $0.9996$ for the Feature branch
and $0.9998$ for the Geo branch. Prediction Spearman correlation was $1.000$
in every Feature and Geo trial, confirming that the controlled reference
expansions did not alter the ordering of the original Hull predictions.

Figure~\ref{fig:hull_geo_stability_four_directions} visualises the geographic
stability result. The same twenty Hull LSOAs remain within the highest-ranked
Geo attribution set under the north, east, south, and west additions, while
none of the newly introduced neighbouring units enters the top-20 set.

% TODO: regenerate from outputs/geo_area_data_eval/n_4000 (area-data ranking)
% before compiling; the existing PDF encodes the coordinates-only ranking.
\begin{figure}[H]
    \centering
    \includegraphics[
        width=0.96\textwidth,
        trim=0.2cm 0.2cm 0.2cm 0.2cm,
        clip
    ]{Chapter4/Figs/4 direction addition.pdf}
    \caption{Geo-SMILE top-20 attribution stability under four cardinal
    neighbouring-LSOA additions. The added unit is outlined in orange.}
    \label{fig:hull_geo_stability_four_directions}
\end{figure}

None of the added LSOAs entered the Geo top 20. Their respective ranks were
82 for the north addition, 57 for east, 101 for south, and 53 for west
(Table~\ref{tab:hull_att_stability_extra}). All four added units have negative
coefficients, meaning that the model depends on their data less than on the
average Hull LSOA. The west neighbour, \texttt{E01013112}, ranked highest of
the four and therefore represented the strongest perturbation of the four
tests. Nevertheless, it ranked only 53rd and remained well outside the
high-attribution set.

\begin{table}[H]
\centering
\caption{Geo-SMILE attribution ranks of the added LSOAs in the cardinal-direction stability experiment.}
\label{tab:hull_att_stability_extra}
\small
\begin{tabular}{llccc}
\toprule
Addition & LSOA code & Extra rank & Extra coef. & In top-20 \\
\midrule
North & \texttt{E01012953} & 82  & $-0.390$ & No \\
East  & \texttt{E01013094} & 57  & $-0.112$ & No \\
South & \texttt{E01013009} & 101 & $-0.694$ & No \\
West  & \texttt{E01013112} & 53  & $-0.065$ & No \\
\bottomrule
\end{tabular}
\end{table}

Figure~\ref{fig:hull_geo_stability_base_vs_west} compares the corresponding
Geo-SMILE rankings for the west cardinal neighbourhood.

Under a fixed inverse kernel, all twenty of the baseline-ranked LSOAs remain in
the top-20 set after the west neighbour is introduced (Jaccard $1.000$,
$\rho=0.9998$). The additional LSOA itself remains outside the
high-attribution set, ranking 53rd. Thus, even under the more restrictive
paired-kernel comparison, the geographic attribution structure is preserved.

% TODO: regenerate from the area-data fixed-kernel west run before compiling.
\begin{figure}[H]
    \centering
    \includegraphics[
        width=0.96\textwidth,
        trim=0.2cm 0.2cm 0.2cm 0.2cm,
        clip
    ]{Chapter4/Figs/stability runs geo base vs west addition version.pdf}
    \caption{Geo-SMILE top-20 ranking before and after the west
    cardinal-neighbour addition under a fixed inverse kernel.}
    \label{fig:hull_geo_stability_base_vs_west}
\end{figure}

The west addition also provides the most conservative Feature illustration,
because it produced the lowest full-vector Feature rank correlation among the
four experiments ($\rho=0.9991$).
Figure~\ref{fig:hull_feature_stability_west} therefore compares the Hull-only
baseline with the west-addition result using the same top-10 Feature layout.
The baseline ranking is led by
\texttt{crime\_rolling\_mean\_6},
\texttt{crime\_rolling\_mean\_3}, and
\texttt{crime\_lag\_1}, followed by socioeconomic, tenure, transport, and
qualification-related predictors. The complete baseline top-10 set is retained
after adding \texttt{E01013112}, giving Jaccard similarity of $1.000$.
The full-vector Spearman correlation of $0.9991$ further indicates that
stability is not confined to top-$k$ membership but extends across the wider
Feature attribution ordering.

The coefficient values displayed in the paired layout are the baseline
$N=4{,}000$ coefficients used to identify the retained top-10 set. The archived
stability output records the west-addition comparison metrics rather than a
separate complete per-feature coefficient vector for the expanded run.


\begin{figure}[H]
    \centering
    {\setlength{\fboxsep}{0pt}%
    \setlength{\fboxrule}{0.6pt}%
    \fbox{%
    \includegraphics[
        width=0.94\textwidth,
        trim=0cm 0cm 0cm 0cm,
        clip
    ]{Chapter4/Figs/feature_stability_west_N4000_corrected.png1.png}}}
    \caption{Feature attribution stability under the west-neighbour addition.
    The top-10 set is fully retained (Jaccard $=1.000$), with near-perfect
    full-vector rank agreement ($\rho=0.9991$).}
    \label{fig:hull_feature_stability_west}
\end{figure}


Taken together, the stability experiments indicate that the Geo-SMILE
explanation is highly insensitive to small changes in the spatial explanation
reference. In the primary four-direction analysis, both Feature and Geo
top-$k$ attribution sets are preserved completely, while the near-unity
full-vector Spearman correlations show that stability extends beyond set
membership to the broader attribution ordering. The supplementary fixed-kernel
west comparison likewise retains all twenty baseline units, with the newly
added neighbour remaining outside the high-attribution set. These results
therefore support the robustness of the Hull Geo-SMILE explanation to modest
changes in spatial reference composition, while avoiding the stronger claim
that attribution values are invariant to all possible spatial perturbations.

\subsubsection{ATT consistency}
\label{sec:hull_att_consistency}

ATT consistency evaluates whether Geo-SMILE produces reproducible attribution coefficients when the same fitted model and the same explanation target are analysed repeatedly. The predictive model, data, surrogate specification, perturbation budget ($N=4{,}000$), and all other experimental settings were held fixed, while independent random seeds were used to generate new perturbation samples for each explanation run. Consequently, variation between runs reflects stochastic variation in the Geo-SMILE explanation procedure.

Following \citet{dehghani_smile_thesis}, consistency was quantified from the run-to-run variance and standard deviation of each attribution coefficient. Lower variance and standard deviation indicate greater consistency because the estimated attribution changes less when the perturbation sample is regenerated.
%For player $j$, let $\beta_j^{(r)}$ denote its coefficient in explanation run
%$r$. The mean attribution coefficient is

%\begin{equation}
%\bar{\beta}_j =
%\frac{1}{R}
%\sum_{r=1}^{R}\beta_j^{(r)},
%\end{equation}

%and its across-run variance and standard deviation are

%\begin{equation}
%s_j^2 =
%\frac{1}{R-1}
%\sum_{r=1}^{R}
%\left(
%\beta_j^{(r)}-\bar{\beta}_j
%\right)^2,
%\qquad
%s_j = \sqrt{s_j^2},
%\end{equation}

%where $R$ denotes the number of independently seeded explanation runs.


Table~\ref{tab:hull_consistency_variance} summarises coefficient variability for the Feature and Geo branches. Feature attributions were highly reproducible across perturbation seeds, with a mean coefficient variance of $0.000227$ and a median variance of $0.000192$. The corresponding mean and median standard deviations were $0.01459$ and $0.01386$, and the largest Feature standard deviation was $0.02396$.

\begin{table}[H]
\centering
\caption{ATT consistency of Geo-SMILE coefficients across independently
seeded perturbation runs}
\label{tab:hull_consistency_variance}
\small
\begin{tabular}{l|rrrrr}
\toprule
Branch &
Mean variance &
Median variance &
Mean SD &
Median SD &
Maximum SD \\
\midrule
Feature &
0.000227 &
0.000192 &
0.01459 &
0.01386 &
0.02396 \\
Geo &
0.01536 &
0.01087 &
0.11045 &
0.10426 &
0.65069 \\
\bottomrule
\end{tabular}
\end{table}

The dominant Feature coefficients were similarly reproducible.
\texttt{crime\_rolling\_mean\_6}, which had the largest absolute mean Feature
coefficient ($-5.2728$), had an across-run SD of only $0.0140$.
Likewise, \texttt{crime\_rolling\_mean\_3} had a mean coefficient of
$-1.3074$ and an SD of $0.0120$. Thus, the large attribution magnitudes of
the principal temporal predictors were preserved across independently
generated perturbation samples.

For the Geo branch, the mean coefficient variance was $0.01536$ and the
median variance was $0.01087$. Mean and median coefficient SDs were
$0.11045$ and $0.10426$, respectively. Most geographic coefficients
therefore exhibited relatively small run-to-run variation. The maximum Geo
SD was $0.65069$ and occurred for LSOA \texttt{E01033104}, which also had by
far the largest mean Geo coefficient ($22.976$). Relative to the magnitude of
that coefficient, the SD is approximately $2.83\%$, indicating that the large
attribution itself is highly reproducible across perturbation seeds.

Across the independently seeded runs, the principal Feature and Geo
attribution magnitudes remained reproducible. The Feature branch showed
very small absolute coefficient variability, while the geographic
attributions remained consistent despite their broader coefficient range.
The persistence of the dominant temporal Feature effects and of the
exceptionally large \texttt{E01033104} Geo effect indicates that the main Geo-SMILE explanation patterns are robust to perturbation
sampling.


\subsection{Additional diagnostics}
\label{sec:hull_additional_metrics}

In addition to the primary evaluation metrics, supplementary diagnostics were
retained to characterise attribution concentration, ranking reproducibility,
and sensitivity to the perturbation budget.

\paragraph{Attribution concentration.}
Using the $N=4{,}000$ metric-suite outputs, Feature attribution remains more
concentrated than Geo attribution. Because the signed Geo coefficient is not a
magnitude, Geo concentration was computed on the non-negative per-LSOA
prediction shift that the coefficient ranks (Figure~\ref{fig:hull_geo_sensitivity}b).
Feature Hoyer sparsity is 0.8491, compared with 0.5161 for Geo; the
corresponding Gini coefficients are 0.7956 and 0.5292, and entropies are
1.7747 and 4.5439. The result indicates that a small group of predictors
accounts for much of the Feature attribution magnitude, whereas geographic
sensitivity is distributed across a broader set of LSOAs despite the single
extreme unit.

\paragraph{Intervention-based attribution impact}
%\label{sec:hull_intervention_impact}

The cross-model ATT Faithfulness analysis evaluates whether attribution quality
tracks predictive competence, but does not directly test whether the
highest-ranked players exert a larger effect on the fitted model than other
players. A complementary intervention experiment was therefore retained for
the final $N=4{,}000$ XGBoost explanation.

For the Feature branch, neutralising the ten highest-ranked factors produced
$W_1=8.3161$, compared with a random-set mean of $4.0700$ and a bottom-ten
response of $0.7733$. The top-minus-random difference was $4.2461$, with an
empirical randomisation $p=0.018$.

For the Geo branch, neutralising all data of the top 20 LSOAs produced
$W_1=37.8354$, compared with a random-set mean of $11.6143$ and a bottom-20
response of $1.4153$. The top-minus-random difference was $26.2211$, with
empirical $p=0.001$, the smallest value attainable with 1,000 random sets. The
Geo top-set response lies well beyond the upper limit of the random-set 95\%
interval ($21.3625$), so the highest-ranked LSOAs are clearly separated from
random spatial sets, while the bottom-ranked LSOAs produce a response far
below the random-set range.

\begin{table}[H]
\centering
\caption{Intervention-based attribution impact for the final
$N=4{,}000$ XGBoost explanation.}
\label{tab:hull_intervention_impact}
\small
\begin{tabular}{lcc}
\toprule
Metric & Feature & Geo \\
\midrule
$k$                  & 10             & 20 \\
Top response         & 8.3161         & 37.8354 \\
Bottom response      & 0.7733         & 1.4153 \\
Random mean          & 4.0700         & 11.6143 \\
Random 95\% interval & 0.8676--8.2044 & 7.1232--21.3625 \\
Top--random          & 4.2461         & 26.2211 \\
Top--bottom          & 7.5428         & 36.4201 \\
Empirical $p$        & 0.0180         & 0.0010 \\
Random sets          & 1,000          & 1,000 \\
\bottomrule
\end{tabular}
\end{table}

This intervention analysis addresses a different property from ATT
Faithfulness. It shows that, for the selected XGBoost explanation, the
highest-ranked players have greater model impact than typical random sets.
It does not establish that attribution quality increases systematically with
predictive-model competence across model families.
\paragraph{Supplementary consistency diagnostics.}
In addition to coefficient variance and standard deviation, rank- and
set-based measures were retained as complementary indicators of explanation
reproducibility. At $N=4{,}000$, mean rank Spearman correlation was 0.831 for
Feature and 0.963 for Geo, while mean score Pearson correlation was 0.9997 and
0.9972, respectively. Mean top-20\% Jaccard agreement was 1.000 for Feature and
0.773 for Geo. These measures describe preservation of attribution ordering,
relative score patterns, and membership of the highest-ranked set, rather than
replacing the coefficient-variability definition of ATT consistency.

\paragraph{Behaviour across perturbation budgets.}
The complete independent sweep is retained as a sensitivity analysis rather
than as a competition in which the largest value on every metric determines
the design. Feature fidelity changes little across the whole range, while Geo
held-out fidelity is highest at $N=4{,}000$. Accuracy fluctuates rather than
increasing monotonically for Feature, with reference overlap between 6/10 and
7/10, whereas Geo overlap is 9/20 at every budget from $N=2{,}000$ upwards and
10/20 at $N=1{,}000$. Faithfulness remains positive throughout and is almost
constant for Geo. Coefficient variability generally decreases as the
perturbation budget increases in both branches, while the supplementary
Geo rank-consistency measure rises monotonically from $0.877$ at $N=1{,}000$ to
$0.983$ at $N=10{,}000$.

\begin{table}[H]
\centering
\caption{Five evaluation diagnostics across independent perturbation budgets
for the Feature and Geo branches.}
\label{tab:hull_metrics_by_n}
\small

\setlength{\tabcolsep}{4pt}

\resizebox{\textwidth}{!}{%
\begin{tabular}{@{}lrcccccc@{}}
\toprule
Branch &
$N$ &
\shortstack{Fidelity\\$R^2$} &
\shortstack{ATT\\accuracy} &
\shortstack{Faithfulness\\$\Delta W_1$} &
\shortstack{Stability\\Jaccard} &
\shortstack{Consistency\\mean variance} &
\shortstack{Consistency\\mean SD} \\
\midrule

\multirow{7}{*}{Feature}
& 1,000  & 0.9764 & 6/10 & 4.229 & 1.000 & 0.000963 & 0.03033 \\
& 2,000  & 0.9774 & 7/10 & 4.280 & 1.000 & 0.000510 & 0.02183 \\
& 3,000  & 0.9768 & 7/10 & 4.262 & 1.000 & 0.000264 & 0.01588 \\
& \textbf{4,000}
& \textbf{0.9770}
& \textbf{6/10}
& \textbf{4.246}
& \textbf{1.000}
& \textbf{0.000227}
& \textbf{0.01459} \\
& 5,000  & 0.9766 & 7/10 & 4.279 & 1.000 & 0.000180 & 0.01305 \\
& 7,500  & 0.9771 & 6/10 & 4.389 & 1.000 & 0.000124 & 0.01082 \\
& 10,000 & 0.9753 & 6/10 & 4.389 & 1.000 & 0.000089 & 0.00914 \\

\midrule

\multirow{7}{*}{Geo}
& 1,000  & 0.9309 & 10/20 & 25.904 & 1.000 & 0.05969 & 0.22625 \\
& 2,000  & 0.9231 & 9/20  & 26.127 & 1.000 & 0.02510 & 0.14931 \\
& 3,000  & 0.9440 & 9/20  & 26.386 & 1.000 & 0.01946 & 0.12320 \\
& \textbf{4,000}
& \textbf{0.9484}
& \textbf{9/20}
& \textbf{26.221}
& \textbf{1.000}
& \textbf{0.01536}
& \textbf{0.11045} \\
& 5,000  & 0.9418 & 9/20  & 26.312 & 1.000 & 0.01814 & 0.11060 \\
& 7,500  & 0.9421 & 9/20  & 26.284 & 1.000 & 0.00831 & 0.08332 \\
& 10,000 & 0.9459 & 9/20  & 26.244 & 1.000 & 0.00584 & 0.06974 \\

\bottomrule
\end{tabular}%
}
\end{table}

Figure~\ref{fig:hull_metrics_by_n_updated} summarises the five Hull evaluation diagnostics across the perturbation-budget sweep.

\begin{figure}[H]
    \centering
    \includegraphics[
        width=0.98\textwidth
    ]{Chapter4/Figs/hull_metrics_n_sweep_area_data.png}
    \caption{Hull evaluation diagnostics across perturbation budgets for the Feature and Geo branches. Stars indicate the final $N=4{,}000$ reporting design.}
    \label{fig:hull_metrics_by_n_updated}
\end{figure}

\subsection{Conclusion}
\label{sec:hull_discussion}

The updated Hull experiment uses a common $N=4{,}000$ perturbation budget for both Feature and Geo explanation, with both branches explaining one fitted XGBoost model. The weighted-Ridge surrogate retains high fidelity in both branches, with held-out $R^2$ of 0.977 for Feature and 0.948 for Geo. The perturbation sweep shows that Feature fidelity is already saturated at lower budgets, whereas Geo fidelity benefits more clearly from moving to 4,000 perturbations.

Recent crime history remains the dominant non-spatial signal. The two rolling crime means have by far the largest attribution magnitudes, and the one-month crime lag is also among the six leading drivers. Socioeconomic and tenure variables contribute at smaller magnitudes. The updated local-map figure therefore replaces the former two-month-lag panel with the one-month-lag panel.

Geographically, the model depends most on the data of a southern and central riverside concentration of LSOAs, led by the extreme unit \texttt{E01033104}. Agreement with the expert reference is substantial for Feature and partial for Geo: the $N=4{,}000$ metric suite recovers 6/10 Feature reference items and 9/20 Geo reference LSOAs, with the Geo agreement concentrated among the most sensitive units. The model-internal diagnostics are strong at \(N=4{,}000\): both branches achieve perfect controlled-input top-\(k\) stability (Jaccard \(=1.000\)), while the consistency analysis shows low run-to-run coefficient variability across independently seeded perturbation samples. The Feature branch has a mean coefficient variance of \(0.000227\) and mean SD of \(0.01459\), whereas the Geo branch has a mean variance of \(0.01536\) and mean SD of \(0.11045\). Because coefficient scales differ between the two branches, these consistency values are interpreted within each branch rather than compared directly across branches.

Faithfulness also remains positive. The top Feature set produces a Wasserstein shift 4.25 above the random-set mean ($p=0.018$), while the top Geo set produces a shift 26.22 above random ($p=0.001$), lying well beyond the upper limit of the random-set distribution.

Taken together, the $N=4{,}000$ results support Geo-SMILE as a reproducible framework for separating non-spatial predictive sensitivity from LSOA-specific geographic sensitivity in the Hull crime-prediction application.

%chapter 5 content should get determined then edit last sentence here
